% BEGIN REV09_TRANSPORTE_POLAR
\subsection{Polar transport of incidence and preservation of readers}
\label{corpus9:iv:polar}

The finite realization of the oriented incidence complex
receives the boundary and orientation constructed in the
APP--TRIT--TPK state. Once the inner products of
this realization are fixed, let \(\mathcal H_0,\mathcal H_1\)
be its finite-dimensional spaces and
\[
 d:\mathcal H_0\longrightarrow\mathcal H_1,\qquad
 \Delta_0=d^*d,\qquad \Delta_1=dd^*.
\]
The incidence and return readers are
\[
 \rho_{\mathrm{int}}(v)=(dv,d^*dv),\qquad
 \rho_{\mathrm{ref}}(a)=(d^*a,dd^*a).
\]
The symbols \(\Delta_0,\Delta_1\) here denote the
Laplacians of this Hilbert-space realization. Orientation,
metrics and the domain of \(d\) are part of its data.

\begin{proposition}[Unitary equivalence of the positive sectors]
\label{corpus9:iv:polar-transporte}
Let
\[
 \mathcal H_0^+=(\ker d)^\perp,\qquad
 \mathcal H_1^+=(\ker d^*)^\perp=\operatorname{Ran}d.
\]
The map
\begin{equation}
 \mathcal U=d(\Delta_0|_{\mathcal H_0^+})^{-1/2}
    :\mathcal H_0^+\longrightarrow\mathcal H_1^+
 \label{corpus9:iv:polar-operador}
\end{equation}
is unitary and satisfies, on those sectors,
\begin{equation}
 \mathcal U\Delta_0=\Delta_1\mathcal U,\qquad
 \mathcal U f(\Delta_0)=f(\Delta_1)\mathcal U
 \label{corpus9:iv:polar-funcional}
\end{equation}
for every function \(f\) defined on the common finite
positive spectrum.
\end{proposition}
\begin{proof}
The reduced singular value decomposition is \(d=V\Sigma W^*\),
where the columns of \(W\) and \(V\) are orthonormal bases
of \(\mathcal H_0^+\) and \(\mathcal H_1^+\), and \(\Sigma\)
is diagonal with strictly positive entries. Hence
\[
 \Delta_0|_{\mathcal H_0^+}=W\Sigma^2W^*,\qquad
 \Delta_1|_{\mathcal H_1^+}=V\Sigma^2V^*,\qquad
 \mathcal U=VW^*.
\]
It follows that \(\mathcal U^*\mathcal U=I_{\mathcal H_0^+}\) and
\(\mathcal U\mathcal U^*=I_{\mathcal H_1^+}\). Evaluating \(f\)
on the same diagonal \(\Sigma^2\) proves
\eqref{corpus9:iv:polar-funcional}. If \(d=0\), both
positive sectors are zero spaces and the identities have
their unique interpretation between zero-dimensional spaces.
\end{proof}

For \(v\in\mathcal H_0^+\) and \(h,t\geq0\), quadratic energy,
the resolvent and the semigroup are preserved:
\begin{align}
 \langle\mathcal Uv,\Delta_1\mathcal Uv\rangle
    &=\langle v,\Delta_0v\rangle,\\
 \mathcal U(I+h\Delta_0)^{-1}
    &=(I+h\Delta_1)^{-1}\mathcal U,\\
 \mathcal Ue^{-t\Delta_0}
    &=e^{-t\Delta_1}\mathcal U.
 \label{corpus9:iv:polar-lectores}
\end{align}
The first identity follows from unitarity and
intertwining. The other two are instances of
\eqref{corpus9:iv:polar-funcional}, with
\(f(x)=(1+hx)^{-1}\) and \(f(x)=e^{-tx}\).

\subsubsection{Realization on the tritic cycle}
In the orthonormal bases of vertices and oriented edges
of the three-element cycle,
\[
 d=
 \begin{pmatrix}
 -1&1&0\\
 0&-1&1\\
 1&0&-1
 \end{pmatrix}.
\]
Matrix multiplication gives
\[
 \Delta_0=\Delta_1=
 \begin{pmatrix}
 2&-1&-1\\
 -1&2&-1\\
 -1&-1&2
 \end{pmatrix}
 =3I-J_3=2I-C-C^{-1},
\]
where every entry of \(J_3\) equals one and \(C\)
is the cyclic permutation. The constant vector spans the
kernel; on its orthogonal complement, \(J_3=0\).
Thus the spectrum is \(\{0,3,3\}\) and
\(\mathcal U=d/\sqrt3\) on the positive sector.
The two readers share this spectrum because they arise from
the same incidence, while retaining their vertex and
edge spaces.

\subsubsection{Orientation and compatibility of refinements}
The replacement \(d\mapsto-d\) preserves \(\Delta_0,\Delta_1\)
and changes \(\mathcal U\) to \(-\mathcal U\). More
generally, for a reversal of edge orientations represented
by a diagonal unitary matrix \(E\),
\[
 d'=Ed,\qquad \Delta_0'=\Delta_0,\qquad
 \Delta_1'=E\Delta_1E^*,\qquad \mathcal U'=E\mathcal U.
\]
These equalities follow by using \(E^*E=I\) in the
definitions. Orientation is thus retained by the
transporter even when the spectral readings
coincide.

In a collection of refinements, the maps
\(J_m^{(0)}:\mathcal H_{0,m}^+\to\mathcal H_{0,m+1}^+\)
and
\(J_m^{(1)}:\mathcal H_{1,m}^+\to\mathcal H_{1,m+1}^+\)
must retain domains and norms and satisfy
\[
 \mathcal U_{m+1}J_m^{(0)}
    =J_m^{(1)}\mathcal U_m
\]
in order to transport this identification across levels.
The proposition proves finite equivalence of the positive
sectors; the preceding square expresses the additional
refinement condition. Kernels, nonlinear constraints
and passages to the limit retain their own proofs.
% END REV09_TRANSPORTE_POLAR
